%% sections/c-policy-memo.tex - H. R. 9510 Bill v4.0 (full bill) - Appendix C
%% Genre: plain-English policy memorandum (TO/FROM/RE/DATE + Q&A). Source
%% deliverable 03. Gray-scale Mermaid figure (TikZ): Figure 8 (verification gate
%% as a system). Single hyphens only; section symbol only in tables.
\billsec{APPENDIX C.}{PLAIN-ENGLISH POLICY MEMORANDUM (NON-OPERATIVE).}\phantomsection\label{toc:aC}

\uscprov{1}{\textit{Independent research draft. Not an enacted law, not pending
legislation, and not legal advice; not endorsed by the FDA, HHS, or any Member of
Congress. This memorandum is non-operative; it explains the bill and carries no
legal force of its own.}}

\uscprov{0}{\textbf{TO:} Health policy staff, Member of the House Committee on Energy
and Commerce (Subcommittee on Health)\\ \textbf{FROM:} CEO Kevin Kawchak,
ChemicalQDevice\\ \textbf{RE:} H. R. 9510, the Verification Before Generation in
Physical AI Oncology Trials Act of 2026 (illustrative number; the Clerk assigns the
real number at introduction), 119th Congress, 2d Session\\ \textbf{DATE:} June 7,
2026}

\uscprov{1}{This memorandum lays out one specific safety gap in how artificial
intelligence is now used in cancer trials, explains why current law does not close
it, and describes a narrow, low-cost bill that would. It is written so a non-engineer
can read it in a few minutes and brief the Member, and it is set out as a short
series of questions and answers.}

\uscprov{1}{\textbf{Q. What is the problem in one sentence?}}
\uscprov{1}{A. Cancer trials are now placing robots in direct physical contact with
patients, what the robot does at the instant it touches a patient is decided by
software that AI increasingly writes in real time, and no Federal law governs the
order of operations, so a sponsor may generate the robot-patient interaction code,
run it on a patient, and document the safety checks afterward. We call this the
sequencing gap.}

\uscprov{1}{\textbf{Q. Why does the order of operations matter so much here?}}
\uscprov{1}{A. Because of embodiment. In an ordinary data pipeline a software error
produces a wrong number a human can catch downstream before anyone is harmed; in an
embodied surgical system the same class of error produces a wrong physical motion a
human cannot catch in time, because one moving body concentrates every failure mode
into a single instrument acting in real time. Checking the code after it has already
moved is, in this setting, checking too late.}

\uscprov{1}{\textbf{Q. Does not existing law already cover this?}}
\uscprov{1}{A. The framework is substantial, and the bill builds on it rather than
replacing it, but every piece stops just short of requiring verification before the
code is generated. The closest analogue is the predetermined change control plan
authority Congress added in 2022 at section 515C of the Federal Food, Drug, and
Cosmetic Act (21 U.S.C. 360e-4), with the FDA final guidance for AI-enabled device
software functions. That authority lets the FDA pre-authorize a planned change before
it ships, which is the right idea, but it is voluntary, submission-based, runs on the
FDA review clock, and is not required in advance of generating new code. The Quality
Management System Regulation (21 CFR part 820, in force since February 2, 2026) and
the part 11 electronic-records rule make a verification record trustworthy but do not
require the verification to come first; the CMS Medicare Advantage guardrail (42 CFR
422.101(c)) keeps a human over AI for a coverage decision but not over robot code;
and Federal nondiscrimination law reaches bias and accessibility but not sequencing
\cite{usc-360e4, fda-pccp-2024, fr-qmsr-2024, cms-4201-ma}.}

\uscprov{1}{A. Several States have acted, but at the claims and utilization-review
layer, not the robot-code layer: California Senate Bill 1120, Illinois House Bill
1806, the Texas Responsible Artificial Intelligence Governance Act (House Bill 149),
and Maryland House Bill 820. No two are identical, and none reaches trial-conduct
code, so a sponsor operating across States cannot build one robot-patient code base
against divergent claims-layer regimes. The substantive content a verification gate
would enforce already exists, scattered across these authorities; the one rule that
makes verification a precondition of generation does not exist anywhere, and that
single missing rule is what H. R. 9510 supplies.}

\uscprov{1}{\textbf{Q. What exactly does the bill do?}}
\uscprov{1}{A. It amends the Federal Food, Drug, and Cosmetic Act (21 U.S.C. 301 et
seq.), current through Public Law 119-93, adding one new section, 515D, codified at 21
U.S.C. 360e-5, placed immediately after the section 515C change control authority it
most resembles. The core rule is a single sequencing requirement: no robot-patient
interaction code may be generated or executed in a Physical AI oncology clinical
investigation unless an automated verification, validation, and uncertainty
quantification process, bound to named external standards, has first cleared for that
interaction, with the record documented and attested. A record created after
generation or execution does not satisfy the section; the timing of the record is
itself the compliance fact. Each check emits ACCEPT, BLOCK, or ESCALATE; the
verification fraction must equal 1.0, with no partial credit; and three gates carry a
hard catastrophe predicate at perfect agreement, tied to ASME V\&V 40, ISO/TS 15066,
and the other named standards so the test is auditable rather than discretionary
\cite{usc-fdca, asmevv40, isots15066}.}

\begin{mermaidfig}
\node[mmperson,text width=26mm] (p) at (0,0) {Person: trial sponsor or investigator};
\node[mmdec,text width=26mm] (gate) at (0,-2.2) {Automated VVUQ gate (new sec. 360e-5)};
\node[mmstep,text width=28mm] (std) at (-3.8,-4.3) {Standards bindings: ASME V\&V 40, ISO/TS 15066};
\node[mmstep,text width=26mm] (rec) at (0,-4.3) {Immutable record (part 11)};
\node[mmlaw,text width=26mm] (fda) at (3.8,-4.3) {FDA review (PMA or PCCP)};
\draw[mmedge] (p) -- node[mmlabel]{submits code} (gate);
\draw[mmedge] (gate) -- (std);
\draw[mmedge] (gate) -- (rec);
\draw[mmedge] (gate) -- node[mmlabel]{ACCEPT} (fda);
\end{mermaidfig}
\vspace{-0.3cm}
\figcaption{Figure 8. The verification gate as a system: a submitted interaction is
checked against named standards, written to an immutable record, and only an ACCEPT
artifact proceeds to FDA review.}

\uscprov{1}{Figure 8 shows the gate as a working system rather than a slogan. A
person, the sponsor or investigator, submits the robot-patient interaction code; the
automated gate under new section 515D checks it against the named external standards
and writes the result to an immutable part 11 record; and only an ACCEPT artifact is
allowed to proceed to FDA review under a premarket application or a predetermined
change control plan. A BLOCK or an ESCALATE never reaches that arrow.}

\uscprov{1}{\textbf{Q. Is it actually feasible to verify code before it runs?}}
\uscprov{1}{A. Yes, and it has already been run on the record, not merely hoped for.
In a first study (VVUQ-01) an automated pipeline held the verification fraction at
1.0 across 51 of 51 tests, accepting 1 candidate artifact, blocking 5, and escalating
1. In a second study (VVUQ-02) a ten-gate assurance suite over the control software
for a surgical humanoid cleared 172 of 172 tests and 32 of 32 sweep iterations, with
a composite mean of 93.56, reproducibly from a fixed seed (20260525). The bill asks
sponsors to do, as a precondition, what has already been demonstrated to work
\cite{kawchak2026vvuq01, kawchak2026vvuq02}.}

\uscprov{1}{\textbf{Q. Why now, and not in a few years?}}
\uscprov{1}{A. Clinical investigation is moving toward real-time, decentralized,
machine-generated control code. A real-time trial cannot wait weeks for a manual,
document-by-document code review between procedures, and manual review already cannot
keep pace with the volume of code a multi-site trial produces in 2026. The gap widens
with every robot added to a site, so an automated gate that clears before generation
is not a brake on modernization; it is the only assurance that can run at the speed
modernization now demands.}

\uscprov{1}{\textbf{Q. How does the bill fit existing law, and what does it cost?}}
\uscprov{1}{A. It is additive, not disruptive. It leaves the investigational-device
exemption, human-subject protection (21 CFR part 50), and the investigational new
drug framework (21 CFR part 312) in place; frames the gate as a pre-authorized change
control discipline of the kind section 515C already contemplates; and routes its
records through the part 11 rule and the part 820 quality system that 2026 sponsors
already meet. A savings clause preserves the State human-in-the-loop laws above. It
creates no new agency or program, directs final regulations within 365 days, and is
expected to be PAYGO neutral, with the statute self-executing in the meantime
\cite{cfr-part50, cfr-part312, cbo-paygo}.}

\uscprov{1}{\textbf{Q. What is the recommendation?}}
\uscprov{1}{A. We recommend that the Member consider introducing H. R. 9510 and refer
it to the Committee on Energy and Commerce (Subcommittee on Health). It closes a
narrow, well-defined, and growing safety gap with a proven, low-cost mechanism; it
builds on tools the FDA and CMS already use; it preserves existing State protections;
and it positions the United States to set the global standard for verifying
surgical-robot software before that software ever touches a patient.}
